\chapter*{Notazione}
\thispagestyle{empty}
\addcontentsline{toc}{chapter}{Notazione}

Questa pagina raccoglie le convenzioni notazionali adottate nel testo.

\section*{Convenzioni generali}

\begin{description}[leftmargin=2.2cm, style=nextline]
  \item[Vettori] Indicati con una freccia: $\vec{F}$, $\vec{v}$, $\vec{r}$, \ldots
  \item[Versori] Vettori di modulo unitario, indicati con un cappuccio: $\hat{x}, \hat{y}, \hat{z}$ per gli assi cartesiani; $\hat{r}, \hat{t}, \hat{n}$ per le coordinate intrinseche o polari.
  \item[Derivate temporali] Notazione di Newton: $\dot{x} \equiv \dfrac{dx}{dt}$ (velocità), $\ddot{x} \equiv \dfrac{d^2x}{dt^2}$ (accelerazione).
  \item[Tensori] Indicati in grassetto, senza freccia: $\mathbf{I}$ per il tensore d'inerzia.
  \item[Pedice $_{CM}$] Grandezza riferita al centro di massa del sistema (es. $\vec{v}_{CM}$).
  \item[Pedice $_0$] Valore iniziale o di riferimento di una grandezza (es. $v_0$, $g_0$).
\end{description}

\section*{Costanti fisiche}

\begin{center}
\begin{tabular}{@{}cll@{}}
\toprule
\textbf{Simbolo} & \textbf{Significato} & \textbf{Valore} \\
\midrule
$g$     & Accelerazione di gravità (superficie terrestre) & $\approx 9{,}81\ \mathrm{m/s^2}$ \\
$G$     & Costante di gravitazione universale              & $6{,}674 \times 10^{-11}\ \mathrm{N\,m^2/kg^2}$ \\
$c$     & Velocità della luce nel vuoto                     & $\approx 3{,}00 \times 10^{8}\ \mathrm{m/s}$ \\
$k_B$   & Costante di Boltzmann                             & $1{,}381 \times 10^{-23}\ \mathrm{J/K}$ \\
$N_A$   & Numero di Avogadro                                & $6{,}022 \times 10^{23}\ \mathrm{mol^{-1}}$ \\
\bottomrule
\end{tabular}
\end{center}

\section*{Simboli dal significato variabile}

Per economia notazionale, alcune lettere --- soprattutto greche --- vengono riutilizzate con significati diversi a seconda dell'argomento, secondo l'uso più comune in quell'ambito della fisica: ad esempio $\gamma$ è il fattore di Lorentz in Relatività Ristretta ma indica un rapporto fra capacità termiche in Termodinamica, $\mu$ è un coefficiente d'attrito in Dinamica ma una massa ridotta negli Urti e in Gravitazione, $\tau$ è un momento meccanico in Dinamica delle Rotazioni ma una costante di tempo negli smorzamenti, $R$ è un raggio in Dinamica delle Rotazioni e in Gravitazione ma la costante dei gas in Termodinamica. Il significato preciso è sempre richiamato nel testo al momento della prima comparsa.
